\documentclass[11pt]{article}

\usepackage[margin=1in]{geometry}
\usepackage{amsmath,amssymb,amsthm}
\usepackage{booktabs}
\usepackage{microtype}
\usepackage[hidelinks,pdftitle={Strong Optimality of Two-Color Permutation Wordle at n=3},pdfauthor={Xiaojian Zeng}]{hyperref}

\newtheorem{theorem}{Theorem}
\newtheorem{lemma}{Lemma}
\newcommand{\bits}{\{0,1\}}
\newcommand{\Om}{\Omega}

\title{Strong Optimality of Two-Color Permutation Wordle at $n=3$}
\author{Xiaojian Zeng\\
  \small PBC School of Finance, Tsinghua University, 100190\\
  \small \texttt{xiaojianzeng@foxmail.com}}
\date{}

\begin{document}
\maketitle

\begin{abstract}
We determine the strong finite-horizon optimum for the two-color Permutation
Wordle game on three ranks.  Among all deterministic adaptive strategies with
legal colored-permutation guesses, the maximum numbers of the 48 secrets solved
within one through six guesses are
\[
  1,\ 8,\ 24,\ 40,\ 47,\ 48.
\]
The operational multicolor \textsc{CircularShift} strategy has exact-round
counts $1,7,16,16,7,1$, and hence attains every one of these optima
simultaneously.  The upper bounds are certified by an unpruned exact Bellman
census over 98,898 state--horizon pairs; a separate verifier reconstructs all
feedback and recomputes all 4,747,104 state--action scores.  This theorem is
restricted to $(n,s)=(3,2)$ and does not settle the general multicolor
conjecture.  We also disclose and resolve, for the theorem, an inconsistency
between the displayed color-update cases in the source definition and the
operational rule supported by its prose, example, and proof.
\end{abstract}

\section{The finite game and its context}

Write $[3]=\{1,2,3\}$ and
\[
  \Om=S_3\times\bits^{[3]}.
\]
An element $(\sigma,\kappa)\in\Om$ consists of a permutation, regarded as a
map from positions to ranks, and a color $\kappa(j)$ attached to each rank
$j$.  Thus there are $6\cdot 2^3=48$ secrets, and the legal-guess alphabet is
the same set.  Against a guess $g=(\gamma,c)$, the feedback is the complete
subset
\[
 F(g;\sigma,\kappa)=
 \{i\in[3]:\gamma(i)=\sigma(i)\text{ and }
 c(\gamma(i))=\kappa(\sigma(i))\}.
 \tag{1}
\]
Full feedback $[3]$ occurs exactly when the guess is the secret.

A deterministic adaptive strategy chooses its next legal guess from the
preceding feedback.  If $T_{\mathcal S}(x)$ is the first round in which it
guesses secret $x$, define
\[
 N_r=\max_{\mathcal S}\#\{x\in\Om:T_{\mathcal S}(x)\le r\}.
 \tag{2}
\]
For a uniform secret this is $48$ times the maximum success probability.
Randomization cannot improve (2): after fixing all random choices, one obtains
a deterministic strategy, and an average cannot exceed the largest
deterministic success count.

Kutin and Smithline introduced Permutation Wordle and its suited, or
multicolor, extension, analyzed \textsc{CircularShift}, and conjectured its
strong optimality at every horizon~\cite{KutinSmithline2026}.  Hiveley studied
the monochrome strategy experimentally~\cite{Hiveley2025}.  Kutin, Smithline,
and Vatter subsequently proved the strong horizon-by-horizon theorem for the
monochrome game, even allowing non-permutation guesses, while retaining the
multicolor analogue as a conjectural direction~\cite{KutinSmithlineVatter2026}.
The result below concerns only the 48-element legal-guess model in (1).

\begin{theorem}\label{thm:main}
For the game in \emph{(1)--(2)},
\[
 (N_1,N_2,N_3,N_4,N_5,N_6)=(1,8,24,40,47,48),
 \tag{3}
\]
and $N_r=48$ for every $r\ge6$.  The operational two-color
\textnormal{\textsc{CircularShift}} strategy defined in Section~\ref{sec:cs} attains $N_r$
for every $r\ge1$.
\end{theorem}

\section{An exact Bellman certificate}\label{sec:bellman}

For a viable-secret set $A\subseteq\Om$, a guess $g\in\Om$, and feedback
$u\subseteq[3]$, let
\[
 A_{g,u}=\{x\in A:F(g;x)=u\}.
\]
Let $V_r(A)$ be the largest number of secrets in $A$ that can be solved in the
next $r$ rounds, and set $V_0(A)=0$.  The following recurrence is both the
computational specification and the upper-bound argument.

\begin{lemma}\label{lem:bellman}
For every $r\ge1$ and $A\subseteq\Om$,
\[
 V_r(A)=\max_{g\in\Om}\left(
    \mathbf 1_{\{g\in A\}}+
    \sum_{u\subsetneq[3]}V_{r-1}(A_{g,u})
 \right).
 \tag{4}
\]
In particular, $N_r=V_r(\Om)$.
\end{lemma}

\begin{proof}
For a fixed first guess $g$, full feedback solves precisely $g$ when $g\in A$.
Each proper feedback $u$ leaves the disjoint viable set $A_{g,u}$, and after
observing $u$ one may choose an optimal continuation for that cell.  This
proves the lower bound in (4).  Conversely, the first guess of any deterministic
strategy is one of the 48 elements $g\in\Om$.  Restricted to the secrets in a
proper response cell, its remaining decision tree solves at most
$V_{r-1}(A_{g,u})$ secrets.  Summing over cells proves the upper bound.
Induction on $r$ completes the proof.
\end{proof}

The exhaustive computation encodes each subset of $\Om$ by an exact 48-bit
mask.  Items are ordered canonically by a lexicographic permutation word and a
three-bit color word.  At every memoized nonempty state the solver evaluates
all 48 legal guesses.  It uses no floating point, symmetry quotient, dominance
rule, or pruning cutoff.  Table~\ref{tab:census} gives the complete census.

\begin{table}[ht]
\centering
\caption{Bellman census and \textsc{CircularShift} distribution.}
\label{tab:census}
\begin{tabular}{rrrrrr}
\toprule
$r$ & state rows & action scores & $N_r$ & exact round & cumulative \\
\midrule
1 & 48,645 & 2,334,960 & 1  & 1  & 1  \\
2 & 37,125 & 1,782,000 & 8  & 7  & 8  \\
3 & 10,743 &   515,664 & 24 & 16 & 24 \\
4 &  2,143 &   102,864 & 40 & 16 & 40 \\
5 &    241 &    11,568 & 47 & 7  & 47 \\
6 &      1 &        48 & 48 & 1  & 48 \\
\midrule
Total & 98,898 & 4,747,104 & & 48 & \\
\bottomrule
\end{tabular}
\end{table}

The released certificate contains all 48 item encodings, all $48^2=2,304$
guess--secret feedback entries, and every state row in the table.  A second
program, which imports no discovery code, reconstructs $\Om$ and (1), checks
every serialized feedback entry, and processes the state rows bottom-up.  For
each row it recomputes the right side of (4) for all 48 guesses, requires every
child dependency to be present, and checks that the stored value is the maximum.
It reports no missing dependency and obtains the six root values in (3).
This discovery/certification separation is also used in other exact finite
computational theorems; for example, Zeng et al. independently re-enumerate a
finite search space and check the decisive certificate from scratch
~\cite{ZengEtAl2026FiniteCertificates}.

Lemma~\ref{lem:bellman} turns this exhaustive integer replay into both a lower
and an upper bound for every deterministic adaptive strategy.  Since the last
root value is $48$, monotonicity and the size of $\Om$ give $N_r=48$ for all
$r\ge6$.

\section{CircularShift attainment and a source discrepancy}\label{sec:cs}

The first guess is the identity permutation with all colors zero.  Following a
nonterminal response, keep every correctly guessed colored rank fixed.  If the
remaining active positions are
$a_1<\cdots<a_\ell$, cyclically move the rank at $a_j$ to $a_{j+1}$ and the
rank at $a_\ell$ to $a_1$.  A moved rank $i$, going clockwise from old position
$p$ to new position $q$ on the cyclic order $1,2,3,1$, has its color incremented
exactly when the arc $(p,q]$ contains its identity starting position $i$.
When $p=q$ (one active position), this is a full circuit and the color is
incremented.  On every viable history in the present two-color instance, the
submitted colors remain in $\bits$ until the secret is solved.

This is the circuit-completion rule described operationally by Kutin and
Smithline and used in their example and proof of the completion-time
formula~\cite[Definition~4.3, Figure~6, and Lemma~4.4]{KutinSmithline2026}.
There is, however, a material inconsistency in that source.  The two displayed
color-update cases in Definition~4.3 do not implement the surrounding
operational description or the case analysis in the proof of Lemma~4.4.  For
the minimal secret encoded as \texttt{213-000}, the first feedback to
\texttt{123-000} is $\{3\}$.  The displayed cyclic condition increments the
color of rank 2 and produces \texttt{213-010}; the circuit rule produces
\texttt{213-000} and solves in round two, as the lemma's formula requires.
Accordingly, Theorem~\ref{thm:main} concerns the precisely stated circuit rule
above and does not assert equivalence with the inconsistent displayed cases.

For a secret $(\sigma,\kappa)$, the operational rule's completion time in the
48-secret census is
\[
 T(\sigma,\kappa)=1+\operatorname{exc}(\sigma)
                    +\sum_{i=1}^3\kappa(i),
 \qquad
 \operatorname{exc}(\sigma)=\#\{i:\sigma(i)>i\}.
 \tag{5}
\]
An independent transcript verifier parses all 168 submitted guesses, checks
their legality, recomputes all 168 feedback sets and all 120 nonsolving
transitions, and verifies (5) for each of the 48 secrets.  Since the excedance
enumerator of $S_3$ is $1+4z+z^2$ and the color-sum enumerator is $(1+z)^3$,
the exact-round generating polynomial is
\[
 (1+4z+z^2)(1+z)^3
 =1+7z+16z^2+16z^3+7z^4+z^5.
 \tag{6}
\]
Thus \textsc{CircularShift} has the exact-round and cumulative columns of
Table~\ref{tab:census}.  Comparing its cumulative counts with the Bellman
optima proves simultaneous attainment and completes the proof of
Theorem~\ref{thm:main}.

\section{Scope, reproducibility, and assistance}

The result is a finite theorem about three ranks, two colors attached to ranks,
and a 48-element legal-guess alphabet.  It says nothing about larger $n$, more
colors, or variants that permit repeated ranks in a guess.  The public sources
state the general multicolor assertion conjecturally, while the later source
also cites multicolor work in preparation~\cite{KutinSmithlineVatter2026}; we
therefore make no priority claim over inaccessible work.

The ancillary material includes the Bellman solver, the independent Bellman
verifier, the complete CSV certificate, the direct \textsc{CircularShift}
simulator, the independent transcript verifier, and their exact JSON summaries.
The decisive computations use integer masks and set operations; no numerical
tolerance or random seed is involved.

OpenAI Codex assisted with literature search, code development, exact
computation, adversarial review organization, and manuscript preparation.
Language-model output is not mathematical evidence.  The proof rests on
Lemma~\ref{lem:bellman}, the complete serialized finite certificate, and the
independent exact replay programs.  No proof assistant was used.

\bibliographystyle{plain}
\bibliography{references}

\end{document}
